\documentclass[12pt,a4paper]{article}

% Chọn XeLaTeX làm compiler trên Overleaf.
\usepackage{fontspec}
\setmainfont{TeX Gyre Termes}
\usepackage[margin=2.5cm]{geometry}
\usepackage{array}
\usepackage{enumitem}

\setlength{\parindent}{0pt}
\setlist[itemize]{leftmargin=1.5em, itemsep=0.4em}

\begin{document}

\begin{center}
    {\Large\textbf{BÁO CÁO TIẾN ĐỘ SPRINT 1}}\\[0.5em]
    \textit{Thời gian thực hiện: 17/08/2026 -- 30/08/2026}
\end{center}

\vspace{1em}

\begin{tabular}{@{}>{\bfseries}p{4cm}p{10cm}@{}}
Nhóm:           & \rule{3cm}{0.4pt} \\[0.6em]
Thành viên:     & Tô Thành Nguyên -- 24521207 \\[0.6em]
Ngày báo cáo:   & 30/08/2026
\end{tabular}

\vspace{1.5em}

\section*{1. Những nội dung đã làm trong 2 tuần qua}

\subsection*{Tô Thành Nguyên -- 24521207}

\begin{itemize}
    \item Khởi tạo dự án Python, yêu cầu phiên bản Python từ 3.11 trở lên và tổ chức package \texttt{laplace}.

    \item Xây dựng lớp cấu hình ứng dụng bằng \texttt{pydantic-settings}, sử dụng tiền tố biến môi trường \texttt{LAPLACE\_}.

    \item Tạo tệp cấu hình mẫu \texttt{.env.example} và bổ sung quy tắc loại trừ các tệp chứa thông tin nhạy cảm.

    \item Tạo entrypoint \texttt{python -m laplace}, hỗ trợ kiểm tra trạng thái môi trường và chạy nhiệm vụ mẫu từ dòng lệnh.

    \item Xây dựng vòng lặp Agent tuần tự cho nhiệm vụ mẫu, có giới hạn số bước để tránh chạy vô hạn.

    \item Bổ sung các trạng thái kết thúc gồm \texttt{completed}, \texttt{failed} và \texttt{step\_limit}.

    \item Bổ sung lý do kết thúc bằng tiếng Việt cho từng trạng thái của Agent.

    \item Kiểm tra dữ liệu đầu vào sai định dạng tại domain boundary và dừng ngay khi một bước thực thi thất bại.

    \item Thiết lập kiểm thử tự động bằng \texttt{pytest} và kiểm tra chất lượng mã nguồn bằng Ruff.

    \item Thiết lập GitHub Actions để chạy kiểm thử trên Python 3.11 và Python 3.12.

    \item Cập nhật README với hướng dẫn cài đặt môi trường, chạy chương trình và kiểm tra chất lượng mã nguồn.

    \item Kết quả tại checkpoint cuối Sprint 1: Ruff đạt, 13 bài kiểm thử đạt, CLI mặc định và CLI nhiệm vụ mẫu chạy thành công.
\end{itemize}

\section*{2. Những nội dung sẽ làm trong 2 tuần tới}

\subsection*{Tô Thành Nguyên -- 24521207}

\begin{itemize}
    \item Hoàn thiện lớp lưu trữ dữ liệu bằng SQLAlchemy 2.0 và SQLite.

    \item Xây dựng các bảng chính gồm người dùng, hội thoại, tin nhắn, tác vụ, bước thực thi, lời gọi công cụ, lời gọi mô hình và trace.

    \item Hoàn thiện nhật ký từng bước của Agent để lưu quyết định, công cụ được sử dụng, kết quả và lý do dừng.

    \item Xây dựng Telegram Bot bằng aiogram 3.x với các lệnh cơ bản, xử lý tin nhắn và tệp đính kèm.

    \item Gắn người dùng Telegram với phiên làm việc riêng và kiểm tra dữ liệu của các người dùng không bị lẫn với nhau.

    \item Xử lý phản hồi dài, trạng thái đang xử lý, cập nhật tiến độ từng bước và yêu cầu hủy lượt đang chạy.

    \item Cài đặt giới hạn tần suất yêu cầu theo từng người dùng.

    \item Kết nối mô hình ngôn ngữ thật và cho phép thay đổi Gemini, OpenAI hoặc Groq bằng cấu hình.

    \item Ghi nhận lượng token, chi phí và độ trễ của từng lời gọi mô hình.

    \item Xây dựng system prompt, schema hành động và kiểm tra tham số công cụ bằng Pydantic.

    \item Thực hiện kiểm thử toàn tuyến từ Telegram đến mô hình ngôn ngữ và phản hồi cho người dùng.
\end{itemize}

\section*{3. Những trở ngại}

\begin{itemize}
    \item Lớp lưu trữ dữ liệu và chức năng trace chưa hoàn tất trong Sprint 1 nên phải chuyển sang đầu Sprint 2.

    \item Cần Telegram Bot Token và API key của nhà cung cấp mô hình để thực hiện kiểm thử với dịch vụ thật.

    \item Các bài kiểm thử hiện tại chủ yếu sử dụng dữ liệu giả lập, chưa thay thế được việc kiểm tra toàn tuyến với Telegram và mô hình thật.

    \item Cần bảo đảm thông tin nhạy cảm trong tệp \texttt{.env} không bị đưa lên repository công khai.
\end{itemize}

\end{document}
